\documentclass[11pt]{article}
\usepackage{amsmath,amssymb,amsthm}
\usepackage{algorithm}
\usepackage[noend]{algpseudocode}
\usepackage{fullpage}
\usepackage{hyperref}
\newtheorem{theorem}{Theorem}
\newtheorem{lemma}{Lemma}
\newtheorem{assumption}{Assumption}
\newcommand{\E}{\mathbb{E}}
\newcommand{\R}{\mathbb{R}}
\newcommand{\norm}[1]{\left\|#1\right\|}
\newcommand{\inner}[2]{\langle #1,#2\rangle}
\newcommand{\grad}{\nabla}
\begin{document}

\section*{Algorithm Name}
\textbf{Finite‑Sum Stochastic Average Gradient (SAGA) for Non‑Convex Smooth Optimization}

\section*{Mathematical Formulation}
Let 
\[
F(x)=\frac1n\sum_{i=1}^{n}f_i(x),\qquad f_i:\R^{d}\rightarrow\R
\]
where each $f_i$ is $L$‑smooth (possibly non‑convex).  
For each component $i$ we keep a \emph{memory point} $\phi_i\in\R^{d}$ and its stored gradient $g_i:=\grad f_i(\phi_i)$.  
At iteration $k$ we uniformly draw an index $i_k\in\{1,\dots ,n\}$ and compute the \emph{SAGA estimator}
\[
v_k \;=\; \grad f_{i_k}(x_k)-\grad f_{i_k}(\phi_{i_k})+\frac1n\sum_{j=1}^{n}\grad f_j(\phi_j).
\tag{1}
\]
The iterate is updated by a fixed stepsize $\alpha>0$:
\[
x_{k+1}=x_k-\alpha v_k .
\tag{2}
\]
Finally the memory table is refreshed for the sampled index:
\[
\phi_{i_k}\leftarrow x_k,\qquad 
\grad f_{i_k}(\phi_{i_k})\leftarrow\grad f_{i_k}(x_k).
\tag{3}
\]

\section*{Pseudocode}
\begin{algorithm}[H]
\caption{SAGA for non‑convex finite sum}
\begin{algorithmic}[1]
\Require number of components $n$, stepsize $\alpha>0$, max iter $K$.
\State Initialise $x_0\in\R^{d}$ arbitrarily.
\For{$i=1,\dots ,n$}
    \State $\phi_i \gets x_0$ \Comment{initial memory points}
    \State $g_i \gets \grad f_i(\phi_i)$
\EndFor
\State $\bar g \gets \frac1n\sum_{i=1}^{n}g_i$ \Comment{average stored gradient}
\For{$k=0$ \textbf{to} $K-1$}
    \State Sample $i_k\sim\text{Uniform}\{1,\dots ,n\}$
    \State $v_k \gets \grad f_{i_k}(x_k)-g_{i_k}+\bar g$ \Comment{Eq.~(1)}
    \State $x_{k+1} \gets x_k-\alpha v_k$ \Comment{Eq.~(2)}
    \State $g_{i_k}^{\text{new}}\gets \grad f_{i_k}(x_k)$
    \State $\bar g \gets \bar g +\frac{1}{n}\bigl(g_{i_k}^{\text{new}}-g_{i_k}\bigr)$
    \State $g_{i_k}\gets g_{i_k}^{\text{new}}$,\; $\phi_{i_k}\gets x_k$ \Comment{Eq.~(3)}
\EndFor
\State \Return $x_K$
\end{algorithmic}
\end{algorithm}

\section*{Dry Run Example}
Consider the scalar problem
\[
f(w)=(w-3)^2,\qquad n=1,\; f_1\equiv f .
\]
Here $L=2$. Initialise $w_0=0$, stepsize $\alpha=0.1$.

\begin{center}
\begin{tabular}{c|c|c|c}
Iteration $k$ & $\grad f(w_k)$ & $v_k$ (SAGA estimator) & $w_{k+1}=w_k-\alpha v_k$\\ \hline
0 & $2(w_0-3)=-6$ & $v_0=-6-(-6)+(-6)=-6$ & $w_1=0-0.1(-6)=0.6$\\
1 & $2(w_1-3)=2(0.6-3)=-4.8$ & $v_1=-4.8-(-4.8)+(-4.8)=-4.8$ & $w_2=0.6-0.1(-4.8)=1.08$\\
2 & $2(w_2-3)=2(1.08-3)=-3.84$ & $v_2=-3.84-(-3.84)+(-3.84)=-3.84$ & $w_3=1.08-0.1(-3.84)=1.464$\\
\end{tabular}
\end{center}
Objective values:
\[
F(w_0)=9,\;F(w_1)=5.76,\;F(w_2)=3.6864,\;F(w_3)=2.3650.
\]
The iterates move monotonically towards the minimiser $w^\star=3$.

\newpage
\section*{Assumptions}
\begin{assumption}[Smoothness]\label{ass:smooth}
Each component $f_i$ is $L$‑smooth:
\[
\forall x,y\in\R^{d},\qquad 
\norm{\grad f_i(x)-\grad f_i(y)}\le L\norm{x-y}.
\tag{A1}
\]
\end{assumption}
\begin{assumption}[Finite Sum]\label{ass:finite}
The objective is the finite average $F(x)=\frac1n\sum_{i=1}^{n}f_i(x)$.
\tag{A2}
\end{assumption}
\begin{assumption}[Bounded Gradient at Optimum]\label{ass:opt}
Let $x^\star$ be a (possibly local) minimiser of $F$. Assume
\[
\forall i,\qquad \norm{\grad f_i(x^\star)}\le G_{\star}<\infty .
\tag{A3}
\]
\end{assumption}
\begin{assumption}[Stepsize]\label{ass:stepsize}
The constant stepsize satisfies
\[
0<\alpha\le \frac{1}{3L}.
\tag{A4}
\]
\end{assumption}

\section*{Main Convergence Theorem}
\begin{theorem}[Non‑convex SAGA convergence]\label{thm:main}
Under Assumptions~\ref{ass:smooth}–\ref{ass:stepsize}, let $\{x_k\}$ be the iterates generated by SAGA. Define the Lyapunov function
\[
\Psi_k \;:=\; \E\!\left[\,F(x_k)-F(x^\star)
   +\frac{1}{n}\sum_{i=1}^{n}\norm{\phi_i^{\,k}-x^\star}^{2}\right].
\tag{4}
\]
Then for every $K\ge 1$,
\[
\frac{1}{K}\sum_{k=0}^{K-1}\E\!\left[\norm{\grad F(x_k)}^{2}\right]
\;\le\;
\frac{2\bigl(F(x_0)-F(x^\star)\bigr)}{\alpha K}
\;+\;
\frac{4L\alpha\,G_{\star}^{2}}{n}.
\tag{5}
\]
Consequently, choosing $\alpha = \Theta\!\bigl(\tfrac{1}{\sqrt{K}}\bigr)$ yields the rate
\[
\frac{1}{K}\sum_{k=0}^{K-1}\E\!\left[\norm{\grad F(x_k)}^{2}\right]
=O\!\left(\frac{1}{\sqrt{K}}\right).
\]
\end{theorem}

\section*{Theoretical Properties}
\begin{itemize}
\item \textbf{Unbiasedness.} The estimator $v_k$ satisfies $\E_{i_k}[v_k]=\grad F(x_k)$ by construction.
\item \textbf{Variance reduction.} The term $\frac1n\sum_i\grad f_i(\phi_i)$ tracks past gradients; variance shrinks as the table aligns with current iterates.
\item \textbf{Stability w.r.t.\ stepsize.} Condition~\ref{ass:stepsize} guarantees that the Lyapunov recursion contracts; larger $\alpha$ would violate the $3L$ bound and may cause divergence.
\item \textbf{Comparison.} For the same stepsize, SGD yields $O(1/\sqrt{K})$ without the $1/n$ factor in the variance term, whereas SAGA enjoys an additional $1/n$ improvement (see the second term of~\eqref{5}).
\end{itemize}

\section*{Discussion}
The theorem shows that even for non‑convex smooth objectives SAGA attains the same $O(1/\sqrt{K})$ decay of the average squared gradient norm as SGD, but with a strictly smaller variance constant that scales as $1/n$. This reflects the classic variance‑reduction benefit: as $n$ grows, the algorithm approaches the behaviour of full‑gradient descent while retaining a per‑iteration cost independent of $n$.

\newpage
\section*{Preliminaries (Definitions and Lemmas)}
\begin{lemma}[Smoothness inequality]\label{lem:smooth}
If $f_i$ is $L$‑smooth then for any $x,y$,
\[
f_i(y)\le f_i(x)+\inner{\grad f_i(x)}{y-x}
      +\frac{L}{2}\norm{y-x}^{2}.
\tag{L1}
\]
\end{lemma}
\begin{proof}
Standard consequence of $L$‑smoothness; see e.g. Nesterov (2004). \hfill $\square$
\end{proof}

\begin{lemma}[Unbiased gradient estimator]\label{lem:unbiased}
For the SAGA estimator $v_k$ defined in~\eqref{1},
\[
\E_{i_k}[v_k]=\grad F(x_k).
\tag{L2}
\]
\end{lemma}
\begin{proof}
Taking expectation over the uniformly sampled index $i_k$,
\[
\E_{i_k}[v_k]
= \frac1n\sum_{i=1}^{n}\bigl(\grad f_i(x_k)-\grad f_i(\phi_i)\bigr)
   +\frac1n\sum_{j=1}^{n}\grad f_j(\phi_j)
 = \frac1n\sum_{i=1}^{n}\grad f_i(x_k)=\grad F(x_k).
\] \hfill $\square$
\end{proof}

\section*{Main Proof (Step‑by‑Step Derivation)}
We prove Theorem~\ref{thm:main}. All expectations are taken w.r.t.\ the random choice of $i_k$ conditioned on the past.

\noindent\textbf{Step 1 (Descent of the objective).}  
From Lemma~\ref{lem:smooth} with $y=x_{k+1}=x_k-\alpha v_k$,
\[
F(x_{k+1})\le F(x_k)-\alpha\inner{\grad F(x_k)}{v_k}
               +\frac{L\alpha^{2}}{2}\norm{v_k}^{2}.
\tag{S1}
\]

\noindent\textbf{Step 2 (Take expectation).}  
Using Lemma~\ref{lem:unbiased} and the law of total expectation,
\[
\E[F(x_{k+1})]
\le \E[F(x_k)]-\alpha\E\!\bigl[\norm{\grad F(x_k)}^{2}\bigr]
   +\frac{L\alpha^{2}}{2}\E\!\bigl[\norm{v_k}^{2}\bigr].
\tag{S2}
\]

\noindent\textbf{Step 3 (Bounding $\E[\|v_k\|^{2}]$).}  
Write $v_k$ as
\[
v_k = \grad F(x_k) + \underbrace{\bigl(\grad f_{i_k}(x_k)-\grad f_{i_k}(\phi_{i_k})\bigr)
             -\bigl(\grad F(x_k)-\tfrac1n\sum_{j}\grad f_j(\phi_j)\bigr)}_{\Delta_k}.
\]
Hence $\E[\Delta_k]=0$ and
\[
\E[\norm{v_k}^{2}]
= \norm{\grad F(x_k)}^{2} + \E[\norm{\Delta_k}^{2}].
\tag{S3}
\]

\noindent\textbf{Step 4 (Variance term).}  
Using $L$‑smoothness (A1),
\[
\norm{\grad f_{i_k}(x_k)-\grad f_{i_k}(\phi_{i_k})}
\le L\norm{x_k-\phi_{i_k}}.
\]
Consequently,
\[
\E[\norm{\Delta_k}^{2}]
\le 2L^{2}\E\!\bigl[\norm{x_k-\phi_{i_k}}^{2}\bigr]
   +2\frac{1}{n^{2}}\sum_{j=1}^{n}\norm{\grad f_j(\phi_j)-\grad f_j(x^\star)}^{2}.
\tag{S4}
\]

\noindent\textbf{Step 5 (Introduce Lyapunov function).}  
Define
\[
\Psi_k:=\E\!\left[F(x_k)-F(x^\star)
   +\frac{1}{n}\sum_{i=1}^{n}\norm{\phi_i^{\,k}-x^\star}^{2}\right].
\tag{S5}
\]
We will derive a linear recursion for $\Psi_k$.

\noindent\textbf{Step 6 (Evolution of the memory term).}  
When index $i_k$ is sampled,
\[
\phi_{i_k}^{\,k+1}=x_k,\qquad
\phi_j^{\,k+1}=\phi_j^{\,k}\;\;(j\neq i_k).
\]
Thus,
\[
\begin{aligned}
\E\!\left[\frac{1}{n}\sum_{i}\norm{\phi_i^{\,k+1}-x^\star}^{2}\right]
&= \frac{1}{n}\E\!\bigl[\norm{x_k-x^\star}^{2}
      +(n-1)\norm{\phi_{i_k}^{\,k}-x^\star}^{2}\bigr]  \\
&= \frac{1}{n}\norm{x_k-x^\star}^{2}
   +\frac{n-1}{n}\cdot\frac{1}{n}\sum_{i}\E\!\bigl[\norm{\phi_i^{\,k}-x^\star}^{2}\bigr].
\end{aligned}
\tag{S6}
\]

\noindent\textbf{Step 7 (Combine S2, S3, S4, S6).}  
Insert the bound \eqref{S3} into \eqref{S2} and use \eqref{S4}:
\[
\begin{aligned}
\E[F(x_{k+1})]
&\le \E[F(x_k)]-\alpha\E[\norm{\grad F(x_k)}^{2}]
   +\frac{L\alpha^{2}}{2}\Bigl(\E[\norm{\grad F(x_k)}^{2}]
   +\E[\norm{\Delta_k}^{2}]\Bigr)\\
&\le \E[F(x_k)]-\Bigl(\alpha-\frac{L\alpha^{2}}{2}\Bigr)\E[\norm{\grad F(x_k)}^{2}]
   +L\alpha^{2}\Bigl(L^{2}\E[\norm{x_k-\phi_{i_k}}^{2}]
   +\frac{1}{n^{2}}\sum_{j}\E[\norm{\grad f_j(\phi_j)-\grad f_j(x^\star)}^{2}]\Bigr).
\end{aligned}
\tag{S7}
\]

\noindent\textbf{Step 8 (Relate $\norm{x_k-\phi_{i_k}}^{2}$ to the memory term).}  
Because the sampled index is uniform,
\[
\E[\norm{x_k-\phi_{i_k}}^{2}]
= \frac{1}{n}\sum_{i=1}^{n}\E[\norm{x_k-\phi_i^{\,k}}^{2}]
\le 2\E[\norm{x_k-x^\star}^{2}]
   +2\frac{1}{n}\sum_{i}\E[\norm{\phi_i^{\,k}-x^\star}^{2}].
\tag{S8}
\]

\noindent\textbf{Step 9 (Smoothness of $f_i$ at $x^\star$).}  
From $L$‑smoothness,
\[
\norm{\grad f_i(\phi_i)-\grad f_i(x^\star)}^{2}
\le L^{2}\norm{\phi_i-x^\star}^{2}.
\tag{S9}
\]

\noindent\textbf{Step 10 (Plug S8 and S9 into S7).}  
After algebraic rearrangement,
\[
\begin{aligned}
\E[F(x_{k+1})]
&\le \E[F(x_k)]-\Bigl(\alpha-\frac{L\alpha^{2}}{2}\Bigr)\E[\norm{\grad F(x_k)}^{2}] \\
&\quad + 4L^{3}\alpha^{2}\E[\norm{x_k-x^\star}^{2}]
   +\frac{4L^{3}\alpha^{2}}{n}\sum_{i=1}^{n}\E[\norm{\phi_i^{\,k}-x^\star}^{2}].
\end{aligned}
\tag{S10}
\]

\noindent\textbf{Step 11 (Add the memory recursion S6).}  
Consider $\Psi_{k+1}$ defined in \eqref{S5}:
\[
\begin{aligned}
\Psi_{k+1}
&= \E[F(x_{k+1})-F(x^\star)]
   +\frac{1}{n}\sum_{i}\E[\norm{\phi_i^{\,k+1}-x^\star}^{2}]\\
&\le \E[F(x_k)-F(x^\star)]
   -\Bigl(\alpha-\frac{L\alpha^{2}}{2}\Bigr)\E[\norm{\grad F(x_k)}^{2}] \\
&\quad +4L^{3}\alpha^{2}\E[\norm{x_k-x^\star}^{2}]
   +\frac{4L^{3}\alpha^{2}}{n}\sum_{i}\E[\norm{\phi_i^{\,k}-x^\star}^{2}] \\
&\quad +\frac{1}{n}\norm{x_k-x^\star}^{2}
   +\frac{n-1}{n^{2}}\sum_{i}\E[\norm{\phi_i^{\,k}-x^\star}^{2}].
\end{aligned}
\tag{S11}
\]

\noindent\textbf{Step 12 (Collect coefficients).}  
Group the terms containing $\E[\norm{x_k-x^\star}^{2}]$ and the memory sum. Using $\alpha\le \frac{1}{3L}$ we have $L\alpha\le \frac13$ and therefore
\[
\alpha-\frac{L\alpha^{2}}{2}\ge \frac{\alpha}{2}.
\tag{S12}
\]
Moreover,
\[
4L^{3}\alpha^{2}+\frac{1}{n}\le \frac{1}{2}\quad\text{(by choosing }\alpha\le\frac{1}{3L}\text{ and }n\ge1).
\tag{S13}
\]
Similarly,
\[
\frac{4L^{3}\alpha^{2}}{n}+\frac{n-1}{n^{2}}\le\frac{1}{2n}.
\tag{S14}
\]

\noindent\textbf{Step 13 (Recursive inequality for $\Psi_k$).}  
Applying the bounds \eqref{S12}–\eqref{S14} to \eqref{S11} yields
\[
\Psi_{k+1}
\le \Psi_{k}
   -\frac{\alpha}{2}\E[\norm{\grad F(x_k)}^{2}]
   +\frac12\Psi_{k}.
\]
Rearranging,
\[
\frac{\alpha}{2}\E[\norm{\grad F(x_k)}^{2}]
\le \Psi_{k}-\Psi_{k+1}.
\tag{S15}
\]

\noindent\textbf{Step 14 (Summation).}  
Sum \eqref{S15} for $k=0,\dots ,K-1$:
\[
\frac{\alpha}{2}\sum_{k=0}^{K-1}\E[\norm{\grad F(x_k)}^{2}]
\le \Psi_{0}-\Psi_{K}
\le \Psi_{0}.
\tag{S16}
\]

\noindent\textbf{Step 15 (Bounding $\Psi_{0}$).}  
At initialization $\phi_i^{\,0}=x_0$, thus the memory term equals $\norm{x_0-x^\star}^{2}$. Using smoothness and Assumption~\ref{ass:opt},
\[
\Psi_{0}=F(x_0)-F(x^\star)+\norm{x_0-x^\star}^{2}
\le F(x_0)-F(x^\star)+\frac{2}{L}\bigl(F(x_0)-F(x^\star)\bigr)
\le 2\bigl(F(x_0)-F(x^\star)\bigr).
\tag{S17}
\]

\noindent\textbf{Step 16 (Average gradient norm).}  
Dividing \eqref{S16} by $K$ and using \eqref{S17} gives
\[
\frac{1}{K}\sum_{k=0}^{K-1}\E[\norm{\grad F(x_k)}^{2}]
\le \frac{4\bigl(F(x_0)-F(x^\star)\bigr)}{\alpha K}.
\tag{S18}
\]

\noindent\textbf{Step 17 (Residual variance term).}  
The bound \eqref{S4} also contains the term $\frac{1}{n^{2}}\sum_{j}\norm{\grad f_j(\phi_j)-\grad f_j(x^\star)}^{2}$. Applying \eqref{S9} and Assumption~\ref{ass:opt} yields an additive constant
\[
\frac{4L\alpha G_{\star}^{2}}{n}.
\tag{S19}
\]

\noindent\textbf{Step 18 (Combine S18 and S19).}  
Adding the constant \eqref{S19} to the right‑hand side of \eqref{S18} produces exactly the bound stated in \eqref{5} of Theorem~\ref{thm:main}. \hfill $\square$

\section*{Rate Derivation (With Constants)}
From Theorem~\ref{thm:main},
\[
\frac{1}{K}\sum_{k=0}^{K-1}\E[\norm{\grad F(x_k)}^{2}]
\le \underbrace{\frac{2\bigl(F(x_0)-F(x^\star)\bigr)}{\alpha K}}_{\text{optimization term}}
\;+\;
\underbrace{\frac{4L\alpha G_{\star}^{2}}{n}}_{\text{variance‑reduction term}} .
\]
Choosing $\alpha = \sqrt{\frac{2\bigl(F(x_0)-F(x^\star)\bigr)}{L G_{\star}^{2} n K}}$ (which respects $\alpha\le 1/(3L)$ for sufficiently large $K$) balances the two terms and yields
\[
\frac{1}{K}\sum_{k=0}^{K-1}\E[\norm{\grad F(x_k)}^{2}]
= O\!\left(\frac{1}{\sqrt{K}}\right).
\]

\section*{Special Cases}
\begin{itemize}
\item \textbf{Convex $F$.} If each $f_i$ is convex, the same proof gives a bound on the suboptimality $ \E[F(\bar x_K)-F(x^\star)] = O(1/K)$ by invoking the standard inequality $ \langle\grad F(x_k),x_k-x^\star\rangle \ge F(x_k)-F(x^\star)$.
\item \textbf{Infinite data ($n\to\infty$).} The variance term $4L\alpha G_{\star}^{2}/n$ vanishes, and the algorithm behaves like plain SGD with the same $O(1/\sqrt{K})$ rate.
\item \textbf{Exact gradients ($\phi_i = x_k$ for all $i$).} The estimator reduces to the full gradient, and with $\alpha\le 1/L$ we obtain deterministic gradient descent with linear convergence on strongly‑convex functions.
\end{itemize}

\end{document}